% !TEX program = lualatex
% Auto-generated from YAML (v3) — 05: デジタル通信の仕組みとセキュリティ

\documentclass[aspectratio=43,professionalfonts,handout]{beamer}
\usepackage{beamer_template}

\newcommand{\LectureNum}{05}
\newcommand{\LectureTitle}{デジタル通信の仕組みとセキュリティ}
\newcommand{\CourseName}{情報リテラシー}
\newcommand{\TextPages}{pp.110-113}
\newcommand{\NextTopic}{情報システムの役割と信頼性}
\newcommand{\NextPages}{pp.114-119}

\begin{document}
% --- Slide 1: title ---

\begin{frame}{\CourseName\quad 第\LectureNum 回「\LectureTitle 」}
  {\small \TermName\quad \TeacherName\quad ／\quad 教科書 \TextPages}

  \vfill

  \begin{block}{今日の流れ}
    \begin{enumerate}
      \item 通信方式（回線交換・パケット交換）
      \item プロトコルとTCP/IP
      \item TCP/IP 4階層モデル
      \item ファイアウォールと暗号化通信
      \item 共通鍵・公開鍵暗号方式
    \end{enumerate}
  \end{block}
\end{frame}

% --- 目標 ---

\begin{frame}{今日の目標}
  \begin{block}{この授業が終わったらできること}
    \begin{itemize}
      \item インターネットの基本的な仕組みを理解し，安全に利用するための注意点を説明できる。
    \end{itemize}
  \end{block}
\end{frame}

% --- Slide 2: twocol ---

\begin{frame}{2つの通信方式}
  \begin{columns}[T]
    \begin{column}{0.48\textwidth}
      \begin{block}{回線交換方式}
        \small
        \begin{itemize}
          \item 通信を始める前に相手との間で通信経路を確立する
          \item いったん経路を確立すれば安定した通信ができる
          \item 回線の障害などで経路が確立できないと通信できない
          \item 例：従来の固定電話
        \end{itemize}
      \end{block}
    \end{column}
    \begin{column}{0.48\textwidth}
      \begin{exampleblock}{パケット交換方式}
        \small
        \begin{itemize}
          \item 情報をパケットとよばれる通信単位に小分けして送信する
          \item あらかじめ経路を確立することなく通信を始められる
          \item パケットが途中で失われたり，混雑の影響を受けたりすることがある
          \item 例：インターネット・電子メール・SNS
        \end{itemize}
      \end{exampleblock}
    \end{column}
  \end{columns}
\end{frame}

% --- Slide 3: terms ---

\begin{frame}{プロトコルとTCP/IP}
  \begin{description}[labelwidth=6em]
    \item[\kw{プロトコル}]
      コンピュータどうしが相互に情報のやりとりをするために，何を使ってどのように送受信するかを決めた約束事
    \item[\kw{TCP/IP（Transmission Control Protocol / Internet Protocol）}]
      インターネットでよく利用されるプロトコル。基本ソフトウェアやアプリケーションに依存せず共通の役割を果たす
    \item[\kw{bps（bits per second）}]
      1秒間に転送できるデータ量の単位。通信速度を表す
  \end{description}
\end{frame}

% --- Slide 4: table ---

\begin{frame}{TCP/IP 4階層モデル}
  \centering
  \begin{tabular}{llll}
    \toprule
    層 & 名称 & 機能 & プロトコルの例 \\
    \midrule
    4層 & アプリケーション層 & アプリケーションに応じた通信の形式や手順を定める & HTTP・SMTP・DNS \\
    3層 & トランスポート層 & パケットが喪失しないで届くことを保証する & TCP・UDP \\
    2層 & インターネット層 & データの行き先と経路を決める & IP \\
    1層 & ネットワークインタフェース層 & 接続されたコンピュータどうしが誤りなくデータ転送する手順を提供する & イーサネット \\
    \bottomrule
  \end{tabular}
\end{frame}

% --- Slide 5: terms ---

\begin{frame}{ファイアウォール}
  \begin{description}[labelwidth=6em]
    \item[\kw{ファイアウォール}]
      ネットワークの出入り口で許されたデータのみを通過させ，それ以外のアクセスを拒絶する仕組み。情報の機密性を守ると同時に，コンピュータウイルスによる障害を防ぎ，可用性を守る
    \item[\kw{パケットフィルタリング方式}]
      パケットのIPアドレスやポート番号などを確認し，通過させるかどうかを判断する方式
    \item[\kw{暗号化}]
      元のデータを第三者に読み取られないよう変換すること。変換前を平文，変換後を暗号文という
    \item[\kw{SSL/TLS}]
      Webサイトとの通信を暗号化する規格。「https://」で始まるサイトで使われる。情報の盗聴・改ざん防止
  \end{description}
\end{frame}

% --- Slide 6: table ---

\begin{frame}{共通鍵暗号方式と公開鍵暗号方式}
  \centering
  \begin{tabular}{lll}
    \toprule
    項目 & 共通鍵暗号方式 & 公開鍵暗号方式 \\
    \midrule
    鍵の種類 & 暗号化・復号に同じ鍵を使う & 暗号化・復号に異なる鍵を使う \\
    処理速度 & 速い & 遅い \\
    鍵の受け渡し & 難しい（安全に共有が必要） & 容易（公開鍵は公開可能） \\
    \bottomrule
  \end{tabular}
\end{frame}

% --- Slide 7: terms ---

\begin{frame}{デジタル署名と認証局}
  \begin{description}[labelwidth=6em]
    \item[\kw{デジタル署名}]
      送信者が秘密鍵で署名し，受信者が公開鍵で検証することで，送信者の正当性とデータの改ざんがないことを確認する仕組み
    \item[\kw{認証局（CA：Certificate Authority）}]
      デジタル証明書を発行し，通信相手の正当性を保証する機関
  \end{description}
\end{frame}

% --- Slide 8: exercise ---

\begin{frame}{演習}
  {\small 各自で取り組んでください。}

  \vfill

  \begin{block}{基本問題}
    \begin{itemize}
      \item 回線交換方式とパケット交換方式の違いを説明せよ。また，災害時にSNSが使いやすい理由をこの違いと関連付けて述べよ
      \item TCP/IPの4階層それぞれの役割を説明せよ
      \item ファイアウォールの役割を説明せよ
    \end{itemize}
  \end{block}

  \vfill

  \begin{exampleblock}{応用問題（余裕がある人）}
    \begin{itemize}
      \item 共通鍵暗号方式と公開鍵暗号方式の違いを説明し，インターネットでの安全な通信（HTTPS）にどちらが使われているか述べよ
    \end{itemize}
  \end{exampleblock}
\end{frame}

% --- Slide 9: assignment ---

\begin{frame}{課題・次回予告}
  \begin{importantbox}[課題（次回までに提出）]
    教科書 pp.110-113 を復習すること。次週は情報システムの役割と信頼性を学ぶ（pp.114-119）
  \end{importantbox}

  \vfill

  \begin{block}{次回}
    「\NextTopic 」（教科書 \NextPages ）
  \end{block}

  \vfill

  {\small
  質問があれば授業後またはTeamsで受け付けます。}
\end{frame}

% --- チェックリスト ---

\begin{frame}{まとめ・チェックリスト}
  \begin{block}{自己チェック（できたら $\checkmark$ をつけよう）}
    \begin{itemize}
      \item[$\square$] 回線交換方式とパケット交換方式の違い・TCP/IP 4階層を理解した
      \item[$\square$] ファイアウォールと暗号化通信（SSL/TLS）の仕組みを学んだ
      \item[$\square$] 共通鍵暗号方式と公開鍵暗号方式の違い・デジタル署名・認証局を把握した
    \end{itemize}
  \end{block}

  \vfill

  {\small 質問があれば授業後またはTeamsで受け付けます。}
\end{frame}

\end{document}
